\section{Introduction}
\label{sec:intro}
Autonomous cyber-physical systems such as cars, drones and industrial robots run a processing pipeline that typically consists of sensing, perception, planning and control. The execution path from a sensor reading to the resulting actuation must complete within an end-to-end deadline for the system to be considered functional, safe, or both. Some steps of the pipeline require hardware acceleration and routinely offload work to GPUs. Such systems are usually built on a middleware, a software layer between the application and the operating system that provides communication and composition, and thereby raises the level of abstraction at which components are written. This is what lets the hundreds of components making up the pipeline be implemented with reasonable effort.

ROS~2~\cite{ros2-standard} is one such middleware, widely used in industry and academia. It abstracts an application into three levels: the callback is the unit of computation; a node owns the interfaces through which its callbacks interact with others; and an executor dispatches the ready callbacks of the nodes assigned to it on one or more threads of an OS process.

This creates the difficulty that hinders analysis, verification and validation of such systems: three timelines run at once, and none of them sees the other two (Fig.~\ref{fig:timelines}). From the point of view of the OS, threads are scheduled onto cores; from that of a ROS~2 executor, ready callbacks are dispatched onto its threads; from that of the GPU, offloaded work is scheduled onto its engines. When a callback runs, and how long it takes, is decided on all three: in Fig.~\ref{fig:timelines}, callback 2 is a single execution to the executor, two stretches on a core to the OS, and four pieces of queued work to the GPU, and the time it spends waiting---for the GPU behind another process's copy, then for a core after being preempted---appears on no single timeline. And the three timelines are only where the work runs; what work there is, and in what order, comes from the application architecture.

Analysing, verifying or validating such a system therefore needs the architecture and all three timelines: which callbacks exist, what triggers each and how they map to executors and threads, when each ran, what GPU work it submitted, and how one's output becomes another's input. Timing analyses of ROS~2~\cite{casini-19-reservationbased,arafat-2026-mtexecutor-cbgroups,teper-2024-end2endmtexecutoranalysisoptimization} take this information as given, and in practice it is rarely written down, or does not reflect the current version; the tools that extract it from a running system each cover the application architecture partially, and none the GPU (Section~\ref{sec:goal}).

\begin{figure}
    \centering
    \includegraphics[width=\columnwidth]{figures/timeline_figure.pdf}
    \caption{Three timelines that emerge as an application runs: the OS scheduling threads (top), a multi-threaded executor dispatching ready callbacks (middle), and the GPU scheduling offloaded work (bottom), where red and green are the work submitted by callback 2. Boxes \textbf{a}--\textbf{c} show the executor's ready set at $t_0-\epsilon$, $t_1$ and $t_6$; callbacks 2 and 3 are mutually exclusive. Callback 1 is ready before $t_0$ but runs only when the OS gives thread 11 a core. Callback 3 is ready from $t_0$ but waits for 2, then for thread 11, until $t_8$. Callback 2 waits for the GPU in $t_2$--$t_3$, its copy queued behind another process's, and for a CPU core in $t_5$--$t_6$, its thread having been preempted. To the executor it is one execution from $t_2$ to $t_7$; to the OS, thread 12 holds a core only in $t_2$--$t_4$ and $t_6$--$t_7$; to the GPU, its work runs in $t_3$--$t_5$. Where callback 2's time went is visible on no single timeline.}
    \label{fig:timelines}
\end{figure}

In this work, we present FISH, a toolset that extracts this information automatically from a traced run of an unmodified ROS~2 application: it is transparent to the application, applies to any application, and records all three timelines on one clock, so that each can be analysed on its own and all of them together. The model is based on actual execution traces. Therefore, it does not require expert knowledge about the application that needs to be modeled, nor does it rely on potentially incomplete or outdated documentation. \textbf{Our contributions are:} \textbf{(I)}~instrumentation extending \texttt{ros2\_tracing}~\cite{bedard-2022-ros2tracing} that discovers the deployed architecture, records every callback execution with its thread---all five kinds the executor dispatches, for rclcpp and rclpy---and attributes GPU work to the callback that submitted it; \textbf{(II)}~an acquisition procedure that traces any launched application unmodified, GPU processes included, and separates its phases; \textbf{(III)}~extraction of a hierarchical model that captures each callback's GPU behaviour per invocation and can be read at any granularity, shown feasible on Autoware.

\section{Challenges and Motivation}
\label{sec:goal}
None of the three timelines can be reconstructed from its own vantage point: the OS sees threads but not what they run, an executor sees callbacks but not where they run, and the GPU sees work but not whose it is. A complete view therefore requires all three, recorded together and tied to one another, and the architecture that decides what runs and in what order. In this section we examine what such a view requires and what stands in the way of constructing it.

\noindent
\textit{\textbf{Application Structure.}}
The application structure is the core of every timeline and necessary to understand functional and timing behaviour: which components the software consists of, which functions are called when those components are triggered, and which channels they use to communicate. Knowing the structure completely means knowing four things. Which executor serves which nodes, and how each executor is configured: single- or multi-threaded, with how many threads, and with which callback groups, which decide which callbacks are allowed to run in parallel. The kind of every dispatchable unit of work---timer, subscription, service, client or waitable---since the executor treats them differently. Every communication pattern between callbacks: producer--consumer over a topic, request--response over a service, sample-and-hold, where a callback reads the latest value another one stored, and join, where a callback fires only once a set of inputs is complete. And the mechanism behind each: a serialised message through the middleware, a pointer handed over inside the process, or state that one callback stores and another later reads, which never becomes a message and can only be inferred.

Neither source code nor current tools capture this structure. Source code does not contain what is decided during execution: which objects a given run creates, whether a delivery stays inside the process, and which callback publishes on which topic---a publisher belongs to the node, not to a callback, and any callback of the node may publish through it. Tools that trace the running system recover parts. ROS-Llama~\cite{blass-2021-rosllama-paper}, Abaza et al.~\cite{abaza-2024-ros2modelsynthesis} and LiME\textsuperscript{IT}~\cite{brandenburg-2026-limeit} recover callbacks and their execution times but assume a single-threaded executor; CARET and its extension~\cite{kuboichi-2022-caret,peng-2022-caretcomp} record nodes, callbacks, topics and executors but not the thread a callback runs on; and \texttt{ros2\_tracing}, on which they build, provides no tracepoints for waitables---the fifth kind the executor dispatches---or for deliveries inside a process, which never reach the client library where its hooks live. Executor configuration is recorded by none of them; and since composable nodes are loaded into containers at run time and handles are reused as objects come and go, even the object list has to be observed rather than read.

\noindent
\textit{\textbf{CPU--GPU Interplay.}}
The CPU and the GPU are two separate domains: they are traced by different tools on different clocks, and what happens on one is invisible from the other. The GPU is moreover a scheduling domain of its own, with several engines---compute and copy---each with its own queue, and a scheduler that orders the submitted work by stream and priority, unseen by both the OS and the executor. Three things follow for a callback that offloads work. It is no longer atomic: what ROS~2 and the OS see as one execution becomes an alternating sequence of CPU segments and accelerator segments~\cite{enright-2024-paam}, and what a CPU-side trace records as its duration mixes computation with waiting for the GPU. Its segments are not fixed either: the GPU work a callback submits branches from invocation to invocation, so the sequence of segments---and hence the timing---differs between invocations of the same callback. And its GPU work must be attributed to it, where the only link between the two domains is the thread the GPU run-time was called from; under a multi-threaded executor that thread is decided at run time, so without the exact thread of every callback execution the callback--GPU relation cannot be established. No existing extractor records the accelerator at all, which leaves the time a callback waits for its GPU work indistinguishable from computation.

\noindent
\textit{\textbf{Task Chains.}}
The executions spread over the three timelines and the two domains are not independent: together they form task chains, from a sensor reading to an actuation, and end-to-end timing can only be computed on a chain. Yet even with everything above in hand---the architecture, the communication and trigger patterns, the mapping of callbacks to nodes and executors---chains are not easily obtained. An application does not do the same thing throughout a run: some communications and executions occur only while it initialises, some only while it shuts down, and the periodic operation in between is what matters from a real-time point of view; some registered callbacks never run at all under a given input. The structure does not carry this distinction; it has to be recovered from what actually executed, and the phases separated before anything in the steady state is measured. Nor does a chain announce itself: in a deployment of hundreds of nodes the graph is entangled, and a chain crosses executors, nodes and processes. Existing tools leave the chain to the user: CARET measures a path the user names, and the extractors above produce a callback graph without the phases or the GPU segments a chain runs through.

What follows from the three challenges is the shape of the solution. The structure supplies the objects; the two domains supply the OS and GPU timelines and the segments of a callback; the chains are what runs across all of them. What is needed is one model that holds all three and the relations of Fig.~\ref{fig:timelines} that tie them together, recorded on one clock and separated by phase, so that any timeline can be reconstructed alone and all of them in combination. Such a model serves more than analysis: it lets a traced run be revisited at any granularity, and it lets the same application be re-examined under a platform the trace did not run on---another core count, scheduler, executor configuration or GPU. All of it has to be obtained without modifying the application, and a model that is observed rather than assumed is only as good as its coverage, so the extraction itself must be checked. Section~\ref{sec:method} describes how FISH builds it.
